\chapter{有理分式}\label{ch:b1:fractions}

\omterm{def:b1:fractions:field}{有理分式}是\omterm{def:b1:poly:def}{多项式}的商。本章虽短，其核心定理
——即部分分式
分解\index{部分分式分解}——把任一
这样的商拆成基本积木 $\frac{c}{(X - a)^k}$ 之和。
除代数上的意义外，它还是积分有理函数
（见\cref{ch:b1:integration}）与求某些级数之和
（见\cref{ch:b1:series}）的标准工具。

\section{域 $K(X)$}

\begin{definition}\label{def:b1:fractions:field}
称 $K$（$= \R$
或 $\C$）上的\emph{有理分式}\index{有理分式}为商 $F = \frac{A}{B}$，其中 $A, B \in K[X]$，$B \neq
0$；两个商 $\frac AB$ 与 $\frac{A'}{B'}$ 在
$AB' = A'B$ 时视为相同。每个分式都有\emph{既约形式}，满足
$\gcd(A, B) = 1$，且在相差常数因子的意义下唯一。在自然的运算下，
有理分式的\omterm{def:b1:logic:sets}{集合} $K(X)$ 构成一个域。

$F$（取既约形式）的\emph{极点}\index{极点}是
$B$ 的根；极点的\emph{阶}是其作为 $B$ 的根的\omterm{def:b1:poly:derivative}{重数}。
$F$ 的\emph{次数}为 $\deg F = \deg A - \deg B \in \Z \cup
\{-\infty\}$。
\end{definition}

\begin{example}[读出极点、阶与次数]\label{ex:b1:fractions:vocab}
设 $F = \dfrac{X^3 - X}{X^4 - 2X^3 + X^2}$，分别分解分子与分母：
分子为 $X(X-1)(X+1)$，分母为 $X^2(X-1)^2$；约去
公因式 $X(X-1)$：
\[
F = \frac{X + 1}{X(X - 1)} \quad\text{（既约形式）.}
\]
\omterm{def:b1:fractions:field}{极点}为 $0$ 与 $1$，二者均为\emph{单}\omterm{def:b1:fractions:field}{极点}——阶应在
既约分母上读出，因此原分母中看似的
二重根无关紧要。次数为 $\deg F = 1 - 2 = -1$，
可由渐近性质看出（$xF(x) \to 1$，当 $x \to \infty$）。
次数的行为与\omterm{def:b1:poly:def}{多项式}的次数相似（$\deg FG = \deg F +
\deg G$，$\deg(F + G) \leq \max$），这一计数规则
在下文求系数时处处用到：每一个“关于
$xF(x)$ 的极限”的论证，实质上都是一次次数计数。
\end{example}

\begin{proposition}[整式部分]\label{prop:b1:fractions:integerpart}
对每个 $F = \frac AB \in K(X)$，都唯一地有 $F = E + \frac RB$，其中 $E
\in K[X]$（即 $F$ 的\emph{整式部分}，或称\omterm{def:b1:poly:def}{多项式}部分），且
$\deg R < \deg B$。$E \neq 0$ 当且仅当 $\deg F \geq 0$。
\end{proposition}

\begin{proof}
由欧几里得除法 $A = BE + R$（见\cref{thm:b1:poly:division}），两边
除以 $B$。唯一性：若 $E + \frac RB = E' + \frac{R'}{B}$，则 $(E -
E')B = R' - R$，且 $\deg(R' - R) < \deg B$，由此推出 $E = E'$，进而 $R
= R'$。
\end{proof}

\begin{example}[先约分，再做除法]\label{ex:b1:fractions:integerpart}
求 $F = \dfrac{X^3 + 1}{X^2 - 1}$ 的整式部分。若不经约分，直接
做除法：$X^3 + 1 = (X^2 - 1)X + (X + 1)$，于是 $F = X +
\frac{X + 1}{X^2 - 1}$。但该分式尚未约分：$X^3 + 1
= (X + 1)(X^2 - X + 1)$ 与 $X^2 - 1 = (X+1)(X-1)$ 有公因式
$X + 1$，且
\[
F = \frac{X^2 - X + 1}{X - 1} = X + \frac{1}{X - 1} :
\]
整式部分仍为 $X$，但真分式部分收缩为一个
基本积木，而 $-1$ 处所谓的“\omterm{def:b1:fractions:field}{极点}”其实从来不是\omterm{def:b1:fractions:field}{极点}。
求\omterm{def:b1:fractions:field}{极点}之前，务必先约分到最简形式：$F$ 的\omterm{def:b1:fractions:field}{极点}
即\emph{既约}分母的根。（整式部分
不受约分影响，这由
\cref{prop:b1:fractions:integerpart}中的唯一性保证。）
\end{example}

\section{$\C$ 上的部分分式分解}

\begin{theorem}[$\C$ 上的分解]\label{thm:b1:fractions:complex}
设 $F = \frac AB \in \C(X)$ 为既约形式，且 $B = c\,(X -
a_1)^{m_1} \cdots (X - a_r)^{m_r}$。则 $F$ 可唯一地写成
\[
F = E + \sum_{i=1}^{r} \sum_{k=1}^{m_i}
\frac{c_{i,k}}{(X - a_i)^k},
\qquad E \in \C[X],\ c_{i,k} \in \C ,
\]
其中 $E$ 为 $F$ 的整式部分。
\end{theorem}

\begin{proof}
由\cref{prop:b1:fractions:integerpart}，不妨设 $\deg A < \deg
B$，并对 $E = 0$ 证明和式分解。

\emph{分离\omterm{def:b1:fractions:field}{极点}。}记 $B = (X - a_1)^{m_1} B_1$，其中
$B_1(a_1) \neq 0$。\omterm{def:b1:poly:def}{多项式} $(X-a_1)^{m_1}$ 与 $B_1$ \omterm{cor:b1:arith:bezout}{互质}
（无公共根），故由 $\C[X]$ 中的贝祖定理（见
\cref{ch:b1:poly}中的注），存在 $U_0, V_0$，使得 $U_0 (X-a_1)^{m_1} +
V_0 B_1 = 1$。两边乘以 $A$，并令 $U = AU_0$，$V =
AV_0$，再除以 $B$：
\[
\frac AB = \frac{V}{(X - a_1)^{m_1}} + \frac{U}{B_1} .
\]
次数条件可以满足：将 $V$ 除以
$(X-a_1)^{m_1}$，设 $V = (X-a_1)^{m_1}Q + V_1$，其中 $\deg V_1 <
m_1$，并把商并入第二项（$U_1 = U +
QB_1$）：
\[
\frac AB = \frac{V_1}{(X - a_1)^{m_1}} + \frac{U_1}{B_1},
\qquad \deg V_1 < m_1 ;
\]
比较次数（$\deg A < \deg B$ 与 $\deg V_1 < m_1$）可推出
$\deg U_1 < \deg B_1$ 同样成立。对 $\frac{U_1}{B_1}$ 迭代，
逐个\omterm{def:b1:fractions:field}{极点}进行，即可把一切化归为下面的单\omterm{def:b1:fractions:field}{极点}情形。

\emph{单\omterm{def:b1:fractions:field}{极点}。}对 $\frac{A}{(X-a)^m}$，且 $\deg A < m$，将 $A$
按 $(X - a)$ 的幂展开，$A = \sum_{j=0}^{m-1} \alpha_j (X - a)^j$
（即\omterm{def:b1:poly:def}{多项式}的泰勒展开，与
\cref{prop:b1:poly:multiplicity}的证明中相同）；作除法即得部分分式
$\frac{\alpha_j}{(X-a)^{m-j}}$。具体地，对 $\frac{X^2 +
1}{(X+2)^3}$，令 $X = Y - 2$，
\[
X^2 + 1 = (Y - 2)^2 + 1 = Y^2 - 4Y + 5 ,
\qquad\text{则}\qquad
\frac{X^2+1}{(X+2)^3} = \frac1{Y} - \frac4{Y^2} + \frac5{Y^3}
\]
其中 $Y = X + 2$：三个部分分式只需对
移位展开式作除法即可得出——遇到单个
高阶\omterm{def:b1:fractions:field}{极点}时最快的做法，亦为推荐用于
\cref{exo:b1:fractions:3}的做法。

\emph{唯一性。}设两种分解相同，其
差为恒等式 $0 = \sum_{i,k} \frac{d_{i,k}}{(X -
a_i)^k}$。两边同乘 $(X - a_1)^{m_1}$，每一项
都带有一个在 $a_1$ 处为零的因子，\emph{除了}
$i = 1$，$k = m_1$ 的那一项，其系数变为 $d_{1,m_1}$ 再加上
至少含一个因子 $(X - a_1)$ 的项。代入
$a_1$ 求值（合法：乘以之后，$a_1$ 处已无
\omterm{def:b1:fractions:field}{极点}），得 $d_{1,m_1} = 0$。最高次项系数既已消去，
以 $(X - a_1)^{m_1 - 1}$ 重复上述过程，并依次向下直至 $k = 1$；
再转向下一个\omterm{def:b1:fractions:field}{极点}。全部 $d_{i,k}$ 皆为零，
故分解唯一。
\end{proof}

\begin{method}[系数的计算]\label{met:b1:fractions:compute}
实际计算中避免使用贝祖方法，应综合运用：
\begin{enumerate}
  \item \emph{最高次幂的遮盖法}：
        $\frac{1}{(X-a)^m}$（$m$ 为\omterm{def:b1:fractions:field}{极点} $a$ 的阶）的系数为
        \[
        c_{a,m} = \Bigl[\,(X - a)^m F\,\Bigr]_{X = a};
        \]
        对 $F = \frac AB$ 的\emph{单}\omterm{def:b1:fractions:field}{极点}，此值为
        $\frac{A(a)}{B'(a)}$;
  \item 在便于计算的点处\emph{求值}，以及
        $xF(x)$ 当 $x \to \infty$ 时的\emph{极限}，以便列出
        其余系数的线性关系；
  \item 若存在，利用\emph{奇偶性}或共轭对称性，以
        减少一半的计算量。
\end{enumerate}
\end{method}

\begin{remark}[部分分式的常见易错点]\label{rem:b1:fractions:pitfalls}
\begin{enumerate}
  \item \emph{忽略整式部分。}部分分式分解
        适用于次数 $< 0$ 的分式，当 $\deg A \geq
        \deg B$ 时，须先作除法
        （见\cref{prop:b1:fractions:integerpart}），否则
        系数的求解会产生矛盾。
  \item \emph{遮盖法超出其适用范围。}
        以 $(X - a)^k$ 乘之并在 $a$ 处求值，只能得到
        $k = m$（即\omterm{def:b1:fractions:field}{极点}的\emph{全阶}）时的系数；
        低阶系数需其他关系（如极限、
        代入求值等）——见\cref{ex:b1:fractions:multiple}。
  \item \emph{忘记约分。}\omterm{def:b1:fractions:field}{极点}应在
        \emph{既约}形式上判断；分子与
        分母间的公因式会产生虚假\omterm{def:b1:fractions:field}{极点}
        （见\cref{ex:b1:fractions:integerpart}）。
  \item \emph{在 $\R$ 上部分分式形式出错。}对不可约
        二次式，分子为\emph{一次式}（$\alpha X +
        \beta$），而非常数；仅写 $\frac{c}{X^2 + 1}$ 会
        丢失解——正确形式由
        \cref{thm:b1:fractions:real}决定，不可随意设定。
\end{enumerate}
\end{remark}

\begin{proof}[单极点公式的证明]
在单\omterm{def:b1:fractions:field}{极点} $a$ 附近：$B = (X - a) Q$，其中 $Q(a) \neq 0$，且 $B' =
Q + (X - a) Q'$，故 $B'(a) = Q(a)$。遮盖值为
$\frac{A(a)}{Q(a)} = \frac{A(a)}{B'(a)}$.
\end{proof}

\begin{example}\label{ex:b1:fractions:simple}
分解 $F = \dfrac{1}{X(X-1)(X-2)}$。共有三个单\omterm{def:b1:fractions:field}{极点}，
逐点使用遮盖法：
\[
c_0 = \frac{1}{(0-1)(0-2)} = \frac12,
\quad
c_1 = \frac{1}{1 \times (1 - 2)} = -1,
\quad
c_2 = \frac{1}{2 \times 1} = \frac12,
\]
于是 $F = \dfrac{1/2}{X} - \dfrac{1}{X-1} + \dfrac{1/2}{X-2}$。
\emph{验算 $X = 3$：}直接算得 $F(3) = \frac{1}{3\cdot2\cdot1} =
\frac16$；由分解式得 $\frac{1/2}{3} - \frac{1}{2} +
\frac{1/2}{1} = \frac16 - \frac12 + \frac12 = \frac16$.
\end{example}

\begin{example}[重极点]\label{ex:b1:fractions:multiple}
分解 $F = \dfrac{X}{(X-1)^2 (X+1)}$。形式为：
$\frac{a}{(X-1)^2} + \frac{b}{X-1} + \frac{c}{X+1}$.
\begin{itemize}
  \item 在二重\omterm{def:b1:fractions:field}{极点}处用遮盖法：$a = \bigl[\frac{X}{X+1}\bigr]_{X=1}
        = \frac12$.
  \item 在 $-1$ 处用遮盖法：$c = \bigl[\frac{X}{(X-1)^2}\bigr]_{X=-1} =
        -\frac14$.
  \item $xF(x)$ 在 $\infty$ 处的极限：$0 = b + c$，故 $b = \frac14$。
\end{itemize}
\[
F = \frac{1/2}{(X-1)^2} + \frac{1/4}{X-1} - \frac{1/4}{X+1} .
\]
\emph{验算 $X = 0$：}$F(0) = 0$，且 $\frac12 - \frac14 - \frac14
= 0$.
\end{example}

\section{$\R$ 上的分解}

\begin{theorem}[$\R$ 上的分解]\label{thm:b1:fractions:real}
设 $F \in \R(X)$ 为既约分式，其分母为
\[
B = c \prod_i (X - a_i)^{m_i} \prod_j (X^2 + p_j X + q_j)^{n_j}
\qquad (p_j^2 - 4q_j < 0).
\]
则 $F$ 可唯一地分解为其整式部分与下列各项之和：
\[
\frac{c_{i,k}}{(X - a_i)^k} \quad (1 \leq k \leq m_i),
\qquad
\frac{\alpha_{j,l}\, X + \beta_{j,l}}{(X^2 + p_j X + q_j)^{l}}
\quad (1 \leq l \leq n_j),
\]
其中系数均为实数。
\end{theorem}

\begin{proof}
在 $\C$ 上分解（见\cref{thm:b1:fractions:complex}）。由于 $F$
为实的，\omterm{def:b1:fractions:field}{极点} $\conj a$ 处（各阶）的系数
是 $a$ 处系数的共轭（对分解式取共轭，
并用唯一性）。将每一对共轭项合并：
\[
\frac{c}{(X - z)^l} + \frac{\conj c}{(X - \conj z)^l}
= \frac{c\,(X - \conj z)^l + \conj c\,(X - z)^l}{\bigl(X^2 - 2\Re(z)X
+ \abs z^2\bigr)^{l}},
\]
其分子等于自身的共轭，因而为实数，次数 $\leq l$；
分出实二次式的倍数后，可将其次数降为
$\leq 1$ 对每一级 $l$ 成立（可用一个小的向下归纳法证明）。实\omterm{def:b1:fractions:field}{极点}的系数
保持实数（共轭运算不改变它们）。唯一性
由 $\C$ 上的唯一性推出。
\end{proof}

\begin{example}[观察共轭配对]\label{ex:b1:fractions:pairing}
以最简单的情形说明证明的机理：在 $\C$ 上，
$\frac1{X^2+1}$ 的\omterm{def:b1:fractions:field}{极点}为 $\pm\iu$，遮盖法所得系数
分别为 $\frac1{2\iu}$（位于 $\iu$ 处）与 $\frac1{-2\iu}$（位于 $-\iu$ 处）——
二者互为共轭，正如定理所预言的：
\[
\frac{1}{X^2 + 1}
= \frac{1/(2\iu)}{X - \iu} - \frac{1/(2\iu)}{X + \iu} .
\]
通分合并：
\[
\frac{1}{2\iu}\cdot\frac{(X + \iu) - (X - \iu)}{X^2 + 1}
= \frac{1}{2\iu}\cdot\frac{2\iu}{X^2 + 1} = \frac{1}{X^2+1} :
\]
虚部相消，实的基本分式完整地重新出现。
对于实被积函数，通常无需离开 $\R$——但当
在复数点处计算\emph{和式}时（如周末问题所做的那样，利用\omterm{def:b1:complex:unity}{单位根}），
复基本分式便是自然的“通货”，
而此配对正是这两种
分解之间的“兑换率”。
\end{example}

\begin{example}\label{ex:b1:fractions:real}
将 $F = \dfrac{4}{(X^2+1)(X-1)^2}$ 在 $\R$ 上分解。形式为
$\frac{aX + b}{X^2 + 1} + \frac{c}{(X-1)^2} + \frac{d}{X - 1}$。
在二重\omterm{def:b1:fractions:field}{极点}处用遮盖法：$c = \bigl[\frac{4}{X^2+1}\bigr]_{X=1} =
2$。在复\omterm{def:b1:fractions:field}{极点} $\iu$ 处（$X^2 + 1$ 对应的分子，
可借助复分解式或直接计算）：两边乘以
$X^2 + 1$ 并令 $X = \iu$：
\[
a\iu + b = \frac{4}{(\iu - 1)^2} = \frac{4}{-2\iu} = 2\iu ,
\]
故 $a = 2$，$b = 0$。$xF(x)$ 在无穷远处的极限给出 $0 = a + d$，故 $d
= -2$。于是
\[
F = \frac{2X}{X^2+1} + \frac{2}{(X-1)^2} - \frac{2}{X-1} .
\]
\emph{验算 $X = 0$：}$F(0) = 4$，且 $0 + 2 + 2 = 4$。
\end{example}

\begin{example}[三次分母的完整计算]\label{ex:b1:fractions:cubic}
将 $F = \dfrac{1}{X^3 + 1}$ 在 $\R$ 上分解。先因式分解：$X^3 +
1 = (X + 1)(X^2 - X + 1)$，二次式的判别式为 $-3 <
0$。形式为 $\frac{a}{X+1} + \frac{bX + c}{X^2 - X + 1}$。遮盖法
在单\omterm{def:b1:fractions:field}{极点} $-1$ 处：$a = \bigl[\frac1{X^2 - X +
1}\bigr]_{X=-1} = \frac13$。$xF(x)$ 在无穷远处的极限给出 $0 = a +
b$，故 $b = -\frac13$。在 $X = 0$ 处求值：$1 = a + c$，故 $c =
\frac23$。于是
\[
\frac{1}{X^3+1} = \frac13\Bigl(\frac{1}{X+1}
+ \frac{-X + 2}{X^2 - X + 1}\Bigr) ,
\]
在 $X = 1$ 处验证：左边为 $\frac12$，右边为
$\frac13\bigl(\frac12 + 1\bigr) = \frac12$。注意此方法之简便：
三个未知数，只需三条简单的线性关系（一次遮盖法、一次
极限、一次求值），无需展开任何式子——
这正是\cref{met:b1:fractions:compute}所述步骤的纯粹形式。
\end{example}

\begin{remark}[用途]
分解之后，有理函数可逐项积分：基本分式
$\frac{1}{(x-a)^k}$ 有初等原函数，而基本分式
$\frac{\alpha x + \beta}{(x^2 + px + q)^l}$ 可化为 $\ln$ 与
$\arctan$（见\cref{ch:b1:integration}）。裂项求和是另一种
常见应用（见\cref{exo:b1:fractions:8}）。
\end{remark}

\begin{remark}[本章的应用]\label{rem:b1:fractions:whereused}
部分分式首先是一个\emph{预处理}步骤：
积分一章（\cref{ch:b1:integration}）将每个有理被积函数
在积分之前都要先经\cref{thm:b1:fractions:real}分解，
而级数一章（\cref{ch:b1:series}）裂项求和有理项时，
方法与\cref{exo:b1:fractions:8}及下文周末问题中完全一致，
——这一技术由此推至极致，
直至 $\sum 1/k^2 = \pi^2/6$。对数导数 $P'/P = \sum
m_i/(X - a_i)$（见\cref{exo:b1:poly:11}）每当研究根的
位置时都会再次出现。超出本册，
$1/\chi(X)$（$\chi$ 为特征\omterm{def:b1:poly:def}{多项式}）的分解
是计算矩阵幂与拉普拉斯变换的基础，这在本科二年级
分册中会用到：基本分式 $\frac{c}{(X - a)^k}$ 正是
在\cref{ch:b1:diffeq}中遇到的解 $t^{k-1}\eu^{at}$ 的代数影子。
\end{remark}

\begin{omfigure}
\begin{tikzpicture}[xscale=1.75, yscale=0.42]
  \draw[->, thin] (-2.7,0) -- (2.8,0) node[below] {$x$};
  \draw[->, thin] (0,-4.6) -- (0,4.9) node[left] {$y$};
  \draw[dashed, gray] (-1,-4.4) -- (-1,4.6);
  \draw[dashed, gray] (1,-4.4) -- (1,4.6);
  \draw[omThm, thin] (-2.6,1.5) -- (2.7,1.5)
    node[above right, font=\small, pos=0.02] {$y = \lambda$};
  \draw[omDef, thick, domain=-2.6:-1.31, samples=80, smooth]
    plot (\x, {1/(\x+1) + 1/\x + 1/(\x-1)});
  \draw[omDef, thick, domain=-0.83:-0.20, samples=120, smooth]
    plot (\x, {1/(\x+1) + 1/\x + 1/(\x-1)});
  \draw[omDef, thick, domain=0.20:0.83, samples=120, smooth]
    plot (\x, {1/(\x+1) + 1/\x + 1/(\x-1)});
  \draw[omDef, thick, domain=1.31:2.7, samples=80, smooth]
    plot (\x, {1/(\x+1) + 1/\x + 1/(\x-1)});
  \fill[omThm] (-0.715,1.5) circle (0.045 and 0.19);
  \fill[omThm] (0.40,1.5) circle (0.045 and 0.19);
  \fill[omThm] (2.3,1.5) circle (0.045 and 0.19);
\end{tikzpicture}

{\small 函数 $F(x) = \frac1{x+1} + \frac1x +
\frac1{x-1}$，属\cref{exo:b1:fractions:12}所述类型，严格
递减于其\omterm{def:b1:fractions:field}{极点} $-1, 0, 1$ 所分隔的各区间上，
（虚线）。每条水平线 $\lambda > 0$ 恰与第一个\omterm{def:b1:fractions:field}{极点}右侧的每个区间相交
一次（圆点标出）：
方程 $F = \lambda$ 的解与\omterm{def:b1:fractions:field}{极点}相互交错。}
\end{omfigure}

\section{练习}

\begin{exercise}[$\star$]\label{exo:b1:fractions:1}
在 $\R$ 上分解：$\dfrac{1}{X^2 - 1}$；
$\;\dfrac{X}{X^2 - 3X + 2}$；$\;\dfrac{X^2 + 1}{X(X-1)}$ \emph{（注意
整式部分）}。
\end{exercise}

\begin{exercise}[$\star$]\label{exo:b1:fractions:2}
在 $\R$ 上分解：$\dfrac{1}{X(X^2 + 1)}$ 与
$\dfrac{X^3}{X^2 + X + 1}$。
\end{exercise}

\begin{exercise}[$\star$]\label{exo:b1:fractions:3}
分解 $\dfrac{1}{X^2(X - 1)}$ 与 $\dfrac{X + 1}{(X - 1)^3}$
\emph{（对第二式，令 $Y = X - 1$）}。
\end{exercise}

\begin{exercise}[$\star\star$]\label{exo:b1:fractions:4}
先在 $\C$ 上、再在 $\R$ 上分解：$\dfrac{1}{X^4 - 1}$。
\end{exercise}

\begin{exercise}[$\star\star$]\label{exo:b1:fractions:5}
在 $\R$ 上分解：$\dfrac{X^2}{(X^2 + 1)^2}$，并由此求
函数 $x \mapsto \dfrac{x^2}{(x^2+1)^2}$ 的一个原函数，已知
$\int \frac{\dd x}{(x^2+1)^2} = \frac12\bigl(\arctan x +
\frac{x}{x^2+1}\bigr) + C$。
\end{exercise}

\begin{exercise}[$\star\star$]\label{exo:b1:fractions:6}
对 $n \in \N^*$，分解 $F_n = \dfrac{n!}{X(X+1)\cdots(X+n)}$
\emph{（单\omterm{def:b1:fractions:field}{极点}位于 $0, -1, \dots, -n$；用遮盖法
并识别出\omterm{def:b1:counting:objects}{二项式系数}）}。
\end{exercise}

\begin{exercise}[$\star\star$]\label{exo:b1:fractions:7}
利用\cref{exo:b1:poly:11}中的恒等式，取 $P = X^n - 1$，
证明下式：
\[
\sum_{k=0}^{n-1} \frac{1}{X - \omega^k} = \frac{n X^{n-1}}{X^n - 1},
\qquad \omega = \eu^{2\iu\pi/n},
\]
并在 $X = 2$、$n = 4$ 时计算两边，以作检验。
\end{exercise}

\begin{exercise}[$\star\star$]\label{exo:b1:fractions:8}
分解 $\dfrac{1}{k(k+1)(k+2)}$，并计算
\[
S_n = \sum_{k=1}^{n} \frac{1}{k(k+1)(k+2)},
\qquad\text{再求}\qquad
\lim_{n \to \infty} S_n .
\]
\end{exercise}

\begin{exercise}[$\star\star\star$]\label{exo:b1:fractions:9}
设 $P \in \R[X]$ 为 $n$ 次\omterm{def:b1:poly:def}{首一多项式}，它有 $n$ 个互异实根
$x_1 < \dots < x_n$。证明下式：
\[
\sum_{i=1}^{n} \frac{1}{P'(x_i)} = 0 \quad (n \geq 2),
\qquad
\sum_{i=1}^{n} \frac{x_i^{\,n-1}}{P'(x_i)} = 1 .
\]
\emph{提示：对 $\frac{X^m}{P}$（$m \leq n - 1$）作分解，并考察
无穷远处系数的衰减——或
利用\omterm{thm:b1:poly:lagrange}{拉格朗日插值}（见\cref{thm:b1:poly:lagrange}），对单项式 $X^m$ 在
节点 $x_i$ 上作插值。}
\end{exercise}

\begin{exercise}[$\star\star$]\label{exo:b1:fractions:10}
将 $\dfrac{1}{X(X+1)^2}$ 在 $\R$ 上分解。承认
$\sum_{k \geq 1} \frac1{k^2} = \frac{\pi^2}6$（已在
本章周末问题中证明），推出
\[
\sum_{n=1}^{\infty} \frac{1}{n(n+1)^2} = 2 - \frac{\pi^2}{6} .
\]
\end{exercise}

\begin{exercise}[$\star\star$]\label{exo:b1:fractions:11}
在 $\R$ 上分解：$\dfrac{1}{(X^2+1)(X^2+4)}$，再分解
$\dfrac{X^2}{(X^2+1)(X^2+4)}$。\emph{提示：两个分母都是
关于 $X^2$ 的\omterm{def:b1:poly:def}{多项式}：先分解 $\frac1{(Y+1)(Y+4)}$。}
\end{exercise}

\begin{exercise}[$\star\star\star$]\label{exo:b1:fractions:12}
（久期方程）设 $F = \sum_{i=1}^{r} \frac{c_i}{X - p_i}$，
$p_1 < p_2 < \dots < p_r$ 为实数，且所有 $c_i > 0$。
\begin{enumerate}
  \item 证明 $F$ 在其定义域的每个区间上
        严格递减，并求 $\pm\infty$ 处的极限，
        以及每个\omterm{def:b1:fractions:field}{极点}两侧的极限。
  \item 由此推出：对每个 $\lambda > 0$，方程 $F(x) =
        \lambda$ 恰有 $r$ 个实根，分别位于
        区间 $\intoo{p_i}{p_{i+1}}$ 内，另有一个位于 $p_r$ 右侧。
        （这类方程支配特征值的摄动；
        此处用到的介值定理在高中水平上使用，
        并在\cref{ch:b1:continuity}中证明。）
\end{enumerate}
\end{exercise}

\section{问题：部分分式的威力}

\begin{problem}\label{pb:b1:fractions:1}
部分分式分解看似只是琐碎的计算，本题
却表明它是一台强大的引擎。以分式
$\frac1{X(X+1)\cdots(X+k)}$ 为输入时，它能以闭式裂项相消地求出整族和式；
以 $\frac1{X^n - 1}$ 为输入时，它给出
三角恒等式，例如
\[
\sum_{k=1}^{n-1}\frac{1}{\sin^2\frac{k\pi}{n}} = \frac{n^2-1}{3} ;
\]
再进一步，这一恒等式推出了数学中最著名的
公式之一，即欧拉的
\[
\sum_{k=1}^{\infty}\frac1{k^2} = \frac{\pi^2}{6}
\]
——这里仅借助本章的代数与
高中三角知识即可得到。\index{裂项相消}\index{巴塞尔问题}
全文记 $\omega = \eu^{2\iu\pi/n}$；数列极限在高中水平上
使用（其形式化见\cref{ch:b1:seq}）。

\medskip
\noindent\textbf{第一部分——裂项相消。}
\begin{enumerate}
  \item 分解 $\frac1{X(X+1)}$，并精确计算 $\sum_{n=1}^{N}
        \frac1{n(n+1)}$，由此可得该和趋于
        $1$。
  \item 对 $\frac1{X(X+2)}$ 作同样处理，证明 $\sum_{n=1}^{N}
        \frac1{n(n+2)} = \frac12\bigl(\frac32 - \frac1{N+1} -
        \frac1{N+2}\bigr) \to \frac34$。（有间隔时，每端各残留\emph{两}个
        边界项。）
  \item 将这一机制形式化：若 $F(X) = G(X) - G(X+1)$，
        其中 $G$ 为在 $\intco1{+\infty}$ 中无\omterm{def:b1:fractions:field}{极点}的某个\omterm{def:b1:fractions:field}{有理分式}，
        则 $\sum_{n=1}^N F(n) = G(1) - G(N+1)$。
        由此求出\cref{exo:b1:fractions:8}中的值 $\frac14$，
        方法是给出 $G$ 作为 $F = \frac1{X(X+1)(X+2)}$ 的见证。
  \item 证明一般的阶乘裂项相消：对 $k \geq 1$，
        \[
        \frac{1}{X(X+1)\cdots(X+k)}
        = \frac1k\biggl(\frac{1}{X(X+1)\cdots(X+k-1)}
        - \frac{1}{(X+1)\cdots(X+k)}\biggr),
        \]
        并由此推出
        \[
        \sum_{n=1}^{\infty}\frac{1}{n(n+1)\cdots(n+k)}
        = \frac{1}{k \cdot k!} .
        \]
        对 $k = 2$ 的情形，与第 3 题核对。
  \item 将\cref{exo:b1:fractions:6}中的分解
        在适当选取的点处求值，证明
        \[
        \sum_{j=0}^{n}(-1)^j\binom nj\,\frac1{j+1}
        = \frac1{n+1},
        \qquad
        \sum_{j=0}^{n}(-1)^j\binom nj\,\frac1{j+2}
        = \frac1{(n+1)(n+2)} .
        \]
\end{enumerate}

\medskip
\noindent\textbf{第二部分——分式 $1/(X^n - 1)$。}
\begin{enumerate}[resume]
  \item 用遮盖法证明：
        \[
        \frac{1}{X^n - 1}
        = \frac1n\sum_{k=0}^{n-1}\frac{\omega^k}{X - \omega^k} .
        \]
  \item 两项检验：对 $n =
        2$ 直接验证公式，并证明 $n$ 个系数之和
        当 $n \geq 2$ 时为零，并说明其必然性（考虑
        $xF(x)$ 在 $x \to \infty$ 时的情形）。
  \item 用遮盖法重新推导
        \cref{exo:b1:fractions:7}中的恒等式：
        $\frac{nX^{n-1}}{X^n-1} = \sum_k \frac1{X - \omega^k}$。
  \item 将共轭\omterm{def:b1:fractions:field}{极点}归并，证明实分解：
        其中 $\theta_k = \frac{2k\pi}n$，
        \[
        \frac{\omega^k}{X - \omega^k} +
        \frac{\omega^{n-k}}{X - \omega^{n-k}}
        = \frac{2\cos\theta_k\,X - 2}
        {X^2 - 2\cos\theta_k\,X + 1} ,
        \]
        并写出
        $\frac1{X^n-1}$ 的完整实分解（区分 $n$ 为奇数与 $n$ 为偶数两种情形）。
  \item 取 $n = 4$ 作特例，并与
        \cref{exo:b1:fractions:4}的结果核对。
\end{enumerate}

\medskip
\noindent\textbf{第三部分——三角和，以及欧拉的
$\pi^2/6$。} 设 $P = 1 + X + \dots + X^{n-1}$，其根为
$\omega, \omega^2, \dots, \omega^{n-1}$（均为单根）。
\begin{enumerate}[resume]
  \item 利用 $\frac{P'}P = \sum_{k=1}^{n-1}\frac1{X - \omega^k}$
        （见\cref{exo:b1:poly:11}），证明
        \[
        \sum_{k=1}^{n-1}\frac{1}{1 - \omega^k} = \frac{n-1}2 .
        \]
  \item 证明 $\dfrac1{1 - \eu^{\iu\theta}} = \dfrac12 +
        \dfrac\iu2\,\frac{\cos(\theta/2)}{\sin(\theta/2)}$，其中
        $\theta \notin 2\pi\Z$（半角分解，
        见\cref{met:b1:complex:trig}），并由第 11 题推出
        $\sum_{k=1}^{n-1}
        \frac{\cos(k\pi/n)}{\sin(k\pi/n)} = 0$——亦可
        由 $k \leftrightarrow n - k$ 的对称性看出。
  \item 对第 11 题的恒等式求导（即利用
        $\bigl(\frac{P'}P\bigr)' = \frac{P''}P -
        \bigl(\frac{P'}P\bigr)^2$ 在 $X = 1$ 处的值），证明
        \[
        \sum_{k=1}^{n-1}\frac{1}{(1 - \omega^k)^2}
        = \frac{(n-1)(5-n)}{12} .
        \]
  \item 记 $\cot t = \frac{\cos t}{\sin t}$，由
        第 12--13 题推出以下两个闭式：
        \[
        \sum_{k=1}^{n-1}\cot^2\frac{k\pi}n = \frac{(n-1)(n-2)}3,
        \qquad
        \sum_{k=1}^{n-1}\frac1{\sin^2\frac{k\pi}n}
        = \frac{n^2-1}3 .
        \]
  \item 对 $n = 3$ 和 $n = 4$，手算验证这两个公式。
  \item 证明不等式 $\cot t < \frac1t < \frac1{\sin t}$
        其中 $t \in \intoo0{\frac\pi2}$（由 $\sin t < t < \tan
        t$），并对 $n = 2m + 1$ 及 $1 \leq k \leq m$ 推出：
        \[
        \cot^2\frac{k\pi}n \;<\; \frac{n^2}{k^2\pi^2}
        \;<\; \frac1{\sin^2\frac{k\pi}n} .
        \]
  \item 对 $k = 1, \dots, m$ 求和这些不等式（利用
        $k \leftrightarrow n - k$ 的对称性，使
        第 14 题的公式减半），并夹逼得到：
        \[
        \sum_{k=1}^{\infty}\frac1{k^2} = \frac{\pi^2}6 .
        \]
\end{enumerate}

\medskip
\noindent\textbf{第四部分——高阶砖块。}
\begin{enumerate}[resume]
  \item 将 $\frac1{X^2-1}$ 的分解式平方，
        并对交叉项重新分解，证明
        \[
        \frac1{(X^2-1)^2}
        = \frac14\Bigl(\frac1{(X-1)^2} + \frac1{(X+1)^2}\Bigr)
        - \frac14\Bigl(\frac1{X-1} - \frac1{X+1}\Bigr),
        \]
        并在 $X = 0$ 处验证。
  \item 结合第 18 题的裂项相消与欧拉的值
        （第 17 题），证明
        \[
        \sum_{n=2}^{\infty}\frac1{(n^2-1)^2}
        = \frac{\pi^2}{12} - \frac{11}{16} ,
        \]
        并数值验证该值精确到小数点后三位。
  \item 对第 8 题的恒等式求导，
        求出 $\sum_{k=0}^{n-1}\frac1{(X -
        \omega^k)^2}$ 的闭形式，并在 $X = 2$、$n = 2$ 时验证。
  \item 证明：对 $k \geq 2$，
        \[
        \sum_{n=k}^{\infty}\frac{1}{\binom nk} = \frac{k}{k-1} .
        \]
        \emph{（将 $1/\binom nk$ 用阶乘写出，
        化归为第 4 题。）}
  \item 取 $k = 3$，给出部分和
        $\sum_{n=3}^{N}\frac1{\binom n3}$ 的精确值及其极限。
\end{enumerate}

\medskip
\noindent\textbf{第五部分——综合。}
\begin{enumerate}[resume]
  \item 作为压轴计算，写出完整的实
        分解 $\dfrac1{X^6 - 1}$，并在 $X =
        0$ 处验证。
  \item 本题究竟在何处用到了：(i) 分解的
        唯一性；(ii) \omterm{def:b1:complex:unity}{单位根}，来自
        \cref{ch:b1:complex}；(iii) 对数导数，来自
        \cref{exo:b1:poly:11}？各用一句话回答。
  \item 综合，用一小段话说明：一个代数恒等式
        ——把分式拆成砖块——生成了精确的
        和、三角恒等式以及 $\pi^2/6$。评述
        代数（精确分解，处处成立）与分析
        （极限、夹逼）之间的分工，并指出
        每条线索在何处被“工业化”：
        裂项相消与比较法在
        \cref{ch:b1:series}中，砖块的积分在
        \cref{ch:b1:integration}中。
\end{enumerate}
\end{problem}
